Data boli stiahnuté z verejného repozitára, ktorý obsahoval surové sekvenačné dáta vo formátu pod5.
Tieto dáta predstavovali čítania DNA reťazcov získané pomocou nanopórového sekvenovania.
Data boli stiahnuté pomocou nástroja aws s použitím príkazu aws s3 cp, ktorý umožňuje kopírovať súbory z Amazon S3 úložiska do lokálneho prostredia.
\begin{lstlisting}[language=bash]
aws s3 sync --no-sign-request s3://ont-open-data/gm24385_2020.09 /media/marek/EAGET/gm24385_2020.09
\end{lstlisting}
Dáta sa stiahli do adresára \texttt{/media/marek/EAGET/gm24385\_2020.09}, a sťahovanie trvalo približne 2 hodiny, pričom celková veľkosť stiahnutých dát bola okolo x GB. %Todo: doplniť presnú veľkosť dát


%popis data a ake su v nich komplikacie
%treba ich prepisat basecallerom dorado do formatu bam
%formate bam ich treba nasledne zoskupit a zarovnat.
Data obsahovali čítania DNA reťazcov, ktoré vznikajú počas procesu nanopórového sekvenovania.
Tieto čítania sú často dlhé a obsahujú chyby, čo predstavuje výzvu pri ich analýze.
Aby bolo možné tieto dáta efektívne analyzovať, bolo potrebné ich prekonvertovať pomocou basecallera Dorado do formátu BAM, ktorý je štandardným formátom pre ukladanie sekvenačných dát a umožňuje efektívne spracovanie a analýzu.
Po konverzii do formátu BAM bolo potrebné dáta zoskupiť a zarovnať, aby sa identifikovali podobné sekvencie a vytvorili konsenzuálne sekvencie, ktoré by mohli reprezentovať pôvodný DNA reťazec.
